\documentclass[10pt,aspectratio=169]{ntuabeamer}

% The framework class already loads the shared NTUA theme, defaults, and macros.
% Add only deck-specific packages here if a copied deck really needs them.
% Bibliography remains deck-controlled, just like in a normal Beamer document.
% If a deck wants .bib support, load biblatex explicitly here.
\usepackage[backend=biber,style=numeric,sorting=none]{biblatex}
% Standard biblatex commands are then available. For example:
% \addbibresource{references.bib}
% Then \slideref{some-bib-key} will use the bib entry if no inline
% \defineslidesource{...}{...} exists for that key.
\addbibresource{references.bib}

% Optional deck-level geometry overrides.
% The framework defaults are intended to work without any tuning.
% Uncomment only when a deck needs non-default logo/header/title spacing.
% \setlength{\titlelogoleftinset}{0.95cm}
% \setlength{\titlelogorightinset}{0.95cm}
% \setlength{\titlelogotopinset}{0.5cm}
% \setlength{\titlelogobottominset}{0.72cm}
% \setlength{\titleblocktopskip}{2.5cm}
% \setlength{\ntuamaintitlelogoheight}{2.0cm}
% \setlength{\schooltitlelogoheight}{1.42cm}
% \setlength{\grouptitlelogoheight}{0.95cm}
% \setlength{\headerlogoheight}{1.02cm}
% \setlength{\headerlogoboxwidth}{1.58cm}
% \setlength{\headercontentheight}{0.86cm}

% Optional logo overrides.
% Uncomment and adapt if this deck should use assets other than the framework defaults.
% \renewcommand{\insertheaderlogo}{%
%   \includegraphics[height=\headerlogoheight]{assets/custom-logo.pdf}%
% }
%
% This starter deck also shows how the title-logo placement hooks can be reused.
% The composition below matches the default layout, but a copied deck can rearrange
% the same building blocks without editing the shared framework.
\renewcommand{\inserttitlelogos}{%
  \begin{tikzpicture}[remember picture,overlay]
    \logotopleft{\insertntuamainlogo}
    \logotopright{\insertschoollogo}
    \logobottomright{\insertgrouplogo}
  \end{tikzpicture}%
}

% Frame-local sources.
% Define reusable citation keys once, then cite them with \slideref{key}.
% Inline definitions are optional convenience. For bibliography files,
% prefer the regular biblatex workflow shown above.
\defineslidesource{goodfellow2014}{Goodfellow et al. \emph{Generative Adversarial Nets}. NeurIPS, 2014.}
\defineslidesource{kingma2015}{Kingma, Ba. \emph{Adam: A Method for Stochastic Optimization}. ICLR, 2015.}
\defineslidesource{vaswani2017}{Vaswani et al. \emph{Attention Is All You Need}. NeurIPS, 2017.}
\defineslidesource{he2016}{He et al. \emph{Deep Residual Learning for Image Recognition}. CVPR, 2016.}
\defineslidesource{lecun1998}{LeCun et al. \emph{Gradient-Based Learning Applied to Document Recognition}. Proceedings of the IEEE, 1998.}

% Title metadata.
\title{NTUA Beamer Example Deck}
\subtitle{Starter examples for layout helpers, logo hooks, and frame-local citations}
\author{NTUA School of Civil Engineering}
\date{\today}

\begin{document}

\begin{frame}[plain]
  \makepresentationtitle
\end{frame}

\begin{frame}{This is an intentionally long example frame title that should wrap to a second line so the shared title bar can be checked visually}
  \begin{itemize}
    \item title wraps to two lines
    \item title and logo should stay vertically centered
    \item use this slide to inspect the shared frame-title layout
  \end{itemize}
\end{frame}

% Basic two-column usage. Optional width arguments let you make the split asymmetric.
\begin{frame}{Two-Column Layout\slideref{goodfellow2014}}
  \twocols[0.53\textwidth][0.43\textwidth]{%
    \textcolor{ntuablue}{Left block}
    \[
      \begin{aligned}
        \min_G \max_D V(D,G)
        &=
        \mathbb{E}_{x \sim p_{\mathrm{data}}}\!\left[\log D(x)\right] \\
        &\quad +
        \mathbb{E}_{z \sim p(z)}\!\left[\log\!\left(1 - D(G(z))\right)\right]
      \end{aligned}
    \]

    This objective follows the original GAN setup\slideref{goodfellow2014}.
  }{%
    \textcolor{ntuablue}{Right block}
    \[
      \mathcal{L}_{\mathrm{gen}} =
      -\mathbb{E}_{z \sim p(z)}\!\left[\log D(G(z))\right]
    \]

    Reusing the same key keeps the same number\slideref{goodfellow2014}.
  }
\end{frame}

% Three-column usage. Widths can also be customized explicitly.
\begin{frame}{Three-Column Layout}
  \threecols[0.24\textwidth][0.30\textwidth][0.38\textwidth]{%
    \textcolor{ntuablue}{SGD}
    \[
      \theta_{t+1} = \theta_t - \eta \nabla_{\theta}\mathcal{L}_t
    \]
  }{%
    \textcolor{ntuablue}{Adam}
    \[
      \theta_{t+1} = \theta_t - \eta \frac{\hat m_t}{\sqrt{\hat v_t} + \epsilon}
    \]
    Adam update\slideref{kingma2015}
  }{%
    \textcolor{ntuablue}{Attention}
    \[
      \begin{aligned}
        S &= \frac{QK^\top}{\sqrt{d_k}} \\
        \operatorname{Attn}(Q,K,V) &= \operatorname{softmax}(S)V
      \end{aligned}
    \]
    Transformer attention\slideref{vaswani2017}
  }
\end{frame}

% Another example of custom-width three-column layout.
\begin{frame}{Custom Three-Column Widths}
  \threecols[0.22\textwidth][0.38\textwidth][0.32\textwidth]{%
    \textcolor{ntuablue}{API}
    \begin{itemize}
      \item \texttt{\textbackslash threecols}
      \item width overrides
    \end{itemize}
  }{%
    \textcolor{ntuablue}{Convolution cost}
    \[
      \operatorname{FLOPs}_{\mathrm{conv}}
      \approx
      2 B H W C_{\mathrm{in}} C_{\mathrm{out}} k^2
    \]
  }{%
    \textcolor{ntuablue}{Memory proxy}
    \[
      \operatorname{Bytes}
      \approx
      4 B H W C
    \]
    A compact example for asymmetric 3-way layouts.
  }
\end{frame}

% Another example of custom-width two-column layout.
\begin{frame}{Custom Width Split}
  \twocols[0.62\textwidth][0.34\textwidth]{%
    \textcolor{ntuablue}{Benchmark formulas}
    \[
      \operatorname{FLOPs}_{\mathrm{dense}} \approx 2 B d_{\mathrm{in}} d_{\mathrm{out}}
    \]
    \[
      \operatorname{Throughput} = \frac{N_{\mathrm{samples}}}{t_{\mathrm{epoch}}}
    \]
    \[
      \operatorname{Speedup} =
      \frac{t_{\mathrm{baseline}}}{t_{\mathrm{candidate}}}
    \]
  }{%
    \textcolor{ntuablue}{Usage}
    \begin{itemize}
      \item custom-width \texttt{\textbackslash twocols}
      \item no citations on this frame
      \item clean footer area
    \end{itemize}
  }
\end{frame}

% Repeated citations on one frame reuse the same number.
% Different keys on the same frame get incrementing numbers in the footer.
\begin{frame}{Per-Frame Citations\slideref{he2016}}
  \small
  \[
    \bd{x}_{\ell+1} = \mathcal{F}(\bd{x}_{\ell}, \bd{W}_{\ell}) + \bd{x}_{\ell}
  \]

  Residual learning\slideref{he2016} appears as citation number 1 on this frame.

  \vspace{0.18cm}

  Adam\slideref{kingma2015} becomes citation number 2 on this same frame.

  \vspace{0.18cm}

  GAN training\slideref{goodfellow2014}, transformer attention\slideref{vaswani2017}, early convolutional benchmarks\slideref{lecun1998}, and a GCN paper pulled from a regular `.bib` file\slideref{kipf2017} add more footer entries.

  \vspace{0.18cm}

  Reusing the residual-learning source again keeps the same mark\slideref{he2016}.
\end{frame}

% Frames with no citations keep a compact footer automatically.
\begin{frame}{Reference-Free Slide}
  \twocols{%
    \textcolor{ntuablue}{Left}
    \[
      \mathcal{J}(\theta) = \frac{1}{N}\sum_{i=1}^{N}
      \ell\!\left(f_{\theta}(x_i), y_i\right)
    \]
  }{%
    \textcolor{ntuablue}{Right}
    \[
      \theta^\star = \arg\min_{\theta} \mathcal{J}(\theta)
    \]
    No citations here, so the footnote area stays empty.
  }
\end{frame}

% Standard biblatex output remains available when a deck wants a classic references slide.
\begin{frame}{Traditional Bibliography}
  \printbibliography[heading=none]
\end{frame}

\end{document}
